
\color{unidarkgreen}{\subsection{Устройство резервирования при отказе выключателя}}
\color{black}

\begin{enumerate}[label=\arabic{section}.\arabic{subsection}.\arabic*,  align=left, labelsep=4pt, leftmargin=0pt, itemindent=57pt]

\item
Устройство резервирования при отказе выключателя (УРОВ) предназначено для ликвидации повреждения в случае отказа выключателя поврежденного присоединения.  

\item
Логика работы УРОВ выполнена с контролем протекания тока через выключатель. Значение по току срабатывания задается параметром \textbf{«Iср УРОВ В»}. Функция УРОВ действует с регулируемой задержкой по времени \textbf{«Тср повт. УРОВ»} на повторное отключение своего выключателя, с задержкой \textbf{«Тср УРОВ»} на отключение смежных выключателей.

\item
Действие УРОВ на смежные выключатели может быть выполнено с контролем сигнала РПВ. Ввод контроля сигнала РПВ в логику работы УРОВ осуществляется программной накладкой \textbf{«УРОВ контр. выкл.»}. Предусмотрен подхват по току, включение которого производится программной накладкой \textbf{«УРОВ подхв.»}.

\item
Алгоритм работы УРОВ ВН приведен на рисунке \ref{pic:UROV}. Алгоритм работы УРОВ СН выполнен аналогично. Уставки УРОВ ВН приведены в таблице \ref{app_set:urovvn}. Уставки УРОВ СН приведены в таблице \ref{app_set:urovsn}.

\vspace{2mm}
\begin{figure}[H]
\centering
\includegraphics[width=1\textwidth,height=1\textheight,keepaspectratio]{\fbpath/053.3341 УРОВ/img/UROV_VN.pdf}
\captionof{figure}{Функциональная схема логики УРОВ ВН}\label{pic:UROV}
\end{figure}

\end{enumerate}